% SGA 3, 第 VI_B 篇, 第 1 轮英文重构底本的中译.
% 范围: 命题 7.8 及完整证明, 包括原编者注 79--80.
% 权威底本: Polo--Gille Exp6B-13oct24.pdf, 组件第 58--59 页
% (合卷第 399--400 页; 重版页眉第 388--389 页; 原版边栏页码第 372--373 页).
% 命题 7.9 始于组件第 60 页, 不在本组件内.
% OCR 与 Reinhold 英文 Markdown 仅为对照材料. 本组件无图.

\medskip
\phantomsection\label{prop:sga3-VIB-7.8}%
\noindent\textbf{命题 7.8.---}
\emph{设 \(G\) 是局部有限型的 \(k\)-群,
\(A\) 与 \(B\) 是 \(G\) 的两个几何约化 \(k\)-群子概形
(因而在其不可约分支的泛点处平滑, 从而由 \(1.3\) 可知平滑).
假设以下条件之一成立:}
\begin{enumerate}
\item[\textup{a)}]\phantomsection\label{cond:sga3-VIB-7.8-a}
  \emph{\(A\) 与 \(B\) 都正规且在 \(k\) 上有限型;}
\item[\textup{b)}]\phantomsection\label{cond:sga3-VIB-7.8-b}
  \emph{\(A\) 连通, 且 \(B\) 在 \(k\) 上有限型.}
\end{enumerate}
\emph{则 \((A,B)\) 在 \(k\) 上有限型,
并表出与 \(A\) 和 \(B\) 在 \(G\) 中的换位子的群预层相伴的
关于 \((\mathrm{fppf})\) 拓扑 (或 \((\mathrm{fpqc})\) 拓扑) 的层.
此外,}%
\footnote{\label{note:sga3-VIB-79}编者注: 以下断言已明确写出,
且在证明中扩充了向代数闭情形的归约.}
\emph{\(k\)-群 \((A,B^0)\) 与 \((A^0,B)\) 都连通, 且}
\phantomsection\label{form:sga3-VIB-7.8-identity-component}%
\[
  (A,B)^0=(A,B^0)\cdot(A^0,B).
\]

\emph{证明.}
由 \hyperref[prop:sga3-VIB-7.6]{7.6},
要证明 \((A,B)\) 是所需的相伴层,
只需证明存在一个整数 \(n\) 使得
\phantomsection\label{form:sga3-VIB-7.8-finite-generation}%
\[
  \nu^n\bigl((A\times_kB)^n\bigr)=(A,B),
\]
这里采用 \hyperref[def:sga3-VIB-7.2.2]{7.2.2} 的记号.
为了证明这一点及其余两个断言, 可以假设 \(k\) 代数闭.

事实上, 设 \(\overline{k}\) 是 \(k\) 的一个代数闭包.
由 \hyperref[concl:sga3-VIB-7.1-iii]{7.1\textup{(iii)}} 及
\hyperref[prop:sga3-VIA-2.4]{VI\(_A\), 2.4},
采用显然的记号, 有
\phantomsection\label{form:sga3-VIB-7.8-base-change}%
\[
\begin{gathered}
  (A,B)_{\overline{k}}
    =(A_{\overline{k}},B_{\overline{k}}),\\
  \bigl((A,B)^0\bigr)_{\overline{k}}
    =(A_{\overline{k}},B_{\overline{k}})^0,\\
  (A,B^0)_{\overline{k}}
    =(A_{\overline{k}},B_{\overline{k}}^0),
  \qquad\text{其余类推.}
\end{gathered}
\]
因此, 假设 \((A_{\overline{k}},B_{\overline{k}})\)
在 \(\overline{k}\) 上有限型;
分别地, 假设 \((A_{\overline{k}},B_{\overline{k}}^0)\)
与 \((A_{\overline{k}}^0,B_{\overline{k}})\) 连通, 且态射\begingroup
\renewcommand{\thefootnote}{\fnsymbol{footnote}}%
\footnote[1]{\label{note:sga3-VIB-7.8-mixed-base-change}
译者注: 法文原文将此乘积的两个因子都印为
\((A_{\overline{k}},B_{\overline{k}}^0)\).
按照所断言的恒等式及整个证明,
第二个因子必须是 \((A_{\overline{k}}^0,B_{\overline{k}})\);
此处恢复了这一混合乘积.}%
\endgroup
\[
  (A_{\overline{k}},B_{\overline{k}}^0)
  \times_{\overline{k}}
  (A_{\overline{k}}^0,B_{\overline{k}})
  \longrightarrow
  (A_{\overline{k}},B_{\overline{k}})^0
\]
是满射. 则由 EGA~IV\(_2\), 2.7.1 和 2.6.1,
在 \(k\) 上相应的断言也成立.

设
\[
  B^1,\ldots,B^p
\]
是 \(B\) 除单位连通分支 \(B^0\) 之外的连通分支.
它们只有有限个, 因为 \(B\) 在 \(k\) 上有限型, 因而是 Noether 型.
在情形 \textup{a)} 中, 同样设
\[
  A^1,\ldots,A^q
\]
是 \(A\) 除 \(A^0\) 之外的连通分支.
在情形 \textup{b)} 中, 只考虑 \(A^0=A\).
设 \(\nu_{ij}\) 为 \(\nu\) 在
\[
  A^i\times_kB^j.
\]
上的限制. 由 \hyperref[cor:sga3-VIA-2.4.1]{VI\(_A\), 2.4.1},
每个 \(A^i\) 与 \(B^j\) 都不可约;
所以 \(A^0\times_kB^j\) 与 \(A^i\times_kB^0\) 也都不可约.
由于 \(G\) 的单位元属于 \(A^0\) 且属于 \(B^0\), 它属于
\[
  \nu_{0j}(A^0\times_kB^j)
  \qquad\text{及}\qquad
  \nu_{i0}(A^i\times_kB^0).
\]
因此, 由 \hyperref[cor:sga3-VIB-7.2.1]{7.2.1},
每个 \(\Gamma_G(\nu_{0j})\) 及每个 \(\Gamma_G(\nu_{i0})\) 都连通.

以 \(u_{0j}\) (分别地, \(u_{i0}\)) 表示
\(\Gamma_G(\nu_{0j})\) (分别地, \(\Gamma_G(\nu_{i0})\))
到 \(G\) 的包含态射. 则
\phantomsection\label{form:sga3-VIB-7.8-generated-connected-groups}%
\[
  (A^0,B)=\Gamma_G\bigl((u_{0j})_{j=0}^{p}\bigr),
  \qquad
  (A,B^0)=\Gamma_G\bigl((u_{i0})_{i=0}^{q}\bigr),
\]
且这些群都连通. 此外, 由 \hyperref[prop:sga3-VIB-7.4]{7.4}
及以上构造可知, 存在一个整数 \(r\),
使得 \((A^0,B)\) 与 \((A,B^0)\) 都包含在
\[
  \nu^r\bigl((A\times_kB)^r\bigr).
\]
中. 在情形 \textup{b)} 中, 有 \((A,B)=(A^0,B)\), 证明已完成.

现在考虑情形 \textup{a)}.%
\footnote{\label{note:sga3-VIB-80}编者注: 以下论证作了扩充,
以纳入新增的\hyperref[cor:sga3-VIB-7.1.1]{推论 7.1.1}.}
已经知道 \((A^0,B)\) 与 \((A,B^0)\) 是 \(G\) 的平滑连通
\(k\)-子群, 因而在 \(k\) 上有限型
(\hyperref[prop:sga3-VIA-2.4]{VI\(_A\), 2.4}).
由于 \(A^0\) 是 \(A\) 的特征子群
(\hyperref[thm:sga3-VIA-2.6.5]{VI\(_A\), 2.6.5}),
它是 \(G\) 的正规子群;
所以 \(\hyperref[concl:sga3-VIB-7.3-v]{7.3\textup{(v)}}\)
表明 \((A^0,B)\) 是 \(G\) 的正规子群,
\((A,B^0)\) 也如此.

因此, 由 \hyperref[cor:sga3-VIB-7.1.1]{7.1.1},
\(G\) 中由 \((A,B^0)\) 与 \((A^0,B)\) 生成的子群 \(H\) 为
\[
  H=(A,B^0)\cdot(A^0,B).
\]
特别地,
\phantomsection\label{form:sga3-VIB-7.8-H-contained}%
\[
  H\subset\nu^{2r}\bigl((A\times_kB)^{2r}\bigr).
\]

若 \(T\) 是 \(G\) 的一个在群法则下稳定的子集 (参见 \(3.0\)),
以 \(T_0\) 表示属于 \(T\) 的 \(G\) 的 \(k\)-值点所成的群.
置
\[
  \Gamma'
    =\bigcup_{q\geq1}\nu^q\bigl((A\times_kB)^q\bigr).
\]
由下文\hyperref[prop:sga3-VIB-7.9]{命题 7.9},
\phantomsection\label{form:sga3-VIB-7.8-rational-commutators}%
\[
  (A^0,B)_0=(A_0^0,B_0),
  \qquad
  (A,B^0)_0=(A_0,B_0^0),
  \qquad
  \Gamma'_0=(A_0,B_0).
\]
所以
\[
  H_0=(A_0^0,B_0)\cdot(A_0,B_0^0)
\]
在 \(\Gamma'_0\) 中的指数有限 (Bible, 第 III 篇, 附录),
因为 \(A_0\) 与 \(B_0\) 正规,
而 \(A_0^0\) (分别地, \(B_0^0\)) 在
\(A_0\) (分别地, \(B_0\)) 中的指数有限.

现在满足了\hyperref[lem:sga3-VIB-7.7]{引理 7.7} 的假设. 由于
\[
  H\subset\nu^{2r}\bigl((A\times_kB)^{2r}\bigr),
\]
存在一个整数 \(m\), 使得
\[
  \nu^{2rm}\bigl((A\times_kB)^{2rm}\bigr)=(A,B),
\]
且 \(G\) 中存在有限多个 \(k\)-值点 \(a_1,\ldots,a_n\), 使得
\[
  (A,B)=a_1H\cup\cdots\cup a_nH.
\]
由于 \(H\) 在 \(k\) 上有限型, 每个 \(a_iH\) 都拟紧.
所以其并 \((A,B)\) 拟紧, 因而在 \(k\) 上有限型.
最后, \(H\) 不可约, 因而每个 \(a_iH\) 都不可约;
且因单位元 \(e\) 属于 \(H\), 可得
\[
  H=(A,B)^0.
\]

% 中文翻译检查点：本组件全部正文、证明、脚注及源注已译；源218行。
% 机械核对：行内数学110处、陈列数学19处及全部引用键守恒（数学中的自然语言已译）。
% 状态：完成源文覆盖和逐段自校；尚未经独立合卷编译与页面审查。
